% CLOCK-GATE REPORT::CLOCKPARITY12::2026-10-18_2112_KST
\documentclass[11pt]{article}
\usepackage{amsmath,amssymb,amsthm,hyperref}
\newtheorem{theorem}{Theorem}
\title{CLOCKPARITY12: Counting-Complexity Firewall and All-Seed NO Trapping}
\date{18 October 2026, 21:12 Asia/Seoul}
\begin{document}\maketitle
\noindent STATUS: PROVED/EXACT restricted theorems (independent review pending), COMPUTATIONAL controls; CLOCK-INDEPENDENT-GATE1 and P versus NP OPEN.

\section*{IDEAS PROPOSED BEFORE LITERATURE}
\item[] A. Test whether checking a supplied witness circuit's full-field polynomial fingerprint is really NP/PH-low; falsifier: a fixed nondegenerate field seed encodes parity-Circuit-SAT.
\item[] B. Combine the prior YES-uncovered lemma CLOCKYESCOVER11 with the dual all-seed-NO-trapped failure; audit target promise quantifiers and selector cost.
\item[] C. Use pairwise-independent collision moments to test whether a source-independent finite random public hash inventory admits even one NO-valid seed for a fixed hard input.
\item[] D. Seek symbolically verifiable witness features rather than assume general circuit-fingerprint evaluation is NP-low.

\section*{CORPUS AND PRIMARY LITERATURE}
Reviewed GitHub main LOSSY-PROMISE-CASCADE6, PARTIAL-SAMPLE-WEDGE7, PROMISE-CASCADE-FIXEDPOINT8, K-ELIMINATION-TRANSFER4, CLOCKSAT1, CLOCKFEATURE5, CLOCKSEED9, CLOCKEXCEPT10, SOURCEQUERY-SYNTHESIS10, MAGRED-STREAM-NO10, PUBLICSEED-PLOTKIN14 and THRES-POLYFIBER10; reviewed prior uncommitted CLOCKYESCOVER11 locally. Searched GitHub for prior supplied-circuit HASHIMAGE parity-verifier results; found no verified equivalent. Literature: Papadimitriou--Zachos 1983, Valiant--Vazirani 1986, Carter--Wegman 1979, CHMY STACS 2021. Previous sampling/Cantelli facts are not claimed novel. Root main.tex not changed because it has no Clock Gate ownership region.

\section*{FULL-DEGREE-PARITY-VERIFIER12}
Let $K=\mathbb F_2[X]/(X^2+X+1)=\mathbb F_4$ and $\alpha=X$, of order 3 and algebraic degree 2. For an explicitly supplied Boolean circuit $C$ on $n$ inputs and $z\in K$ define
\[
\operatorname{FPEQ}_{\alpha}=\{(C,z):\sum_{i=0}^{2^n-1}C(i)\alpha^i=z\}.
\]
\begin{theorem}
$\operatorname{FPEQ}_{\alpha}$ is $\oplus P$-complete under polynomial-time many-one reductions, already for target value $z=1$.
\end{theorem}
\noindent Proof. For any $k$-variable circuit $C$, construct a $(k+2)$-variable circuit $C'$ on integer input $y<2^{k+2}$ which outputs $C(y/3)$ if $3\mid y$ and $y/3<2^k$, and otherwise zero. Division by constant 3, comparison and $C$ have polynomial circuit size. Since $\alpha^3=1$,
\[
\sum_y C'(y)\alpha^y=\sum_{j=0}^{2^k-1}C(j)(\alpha^3)^j=\#\operatorname{SAT}(C)\bmod2.
\]
Thus parity-Circuit-SAT reduces to $\operatorname{FPEQ}_\alpha$ with $z=1$. For membership, each of the two field-coordinate bits is a parity of the accepting assignments weighted by an efficiently computable coefficient of $\alpha^i$; prescribed parity equalities, their complements, and finite conjunctions lie in $\oplus P$ by addition/product of accepting-path counts modulo 2. QED.

\noindent CRUCIAL GUARD: This applies to equality checking for an explicitly supplied candidate $C$, not to the existential language $\exists C:|C|\le t$ and $H_\alpha(TT(C))=z$. In fact the coefficients $1,\alpha$ at positions 0 and 1 span all four labels of $\mathbb F_4$: for $t\ge3n$ every field label has an easy weight-two truth-table representative, so the existential image language is all YES. The verifier can be $\oplus P$-hard while that language is trivial. A direct evaluator enumerating all $N=2^n$ truth-table positions takes $N^{1-o(1)}$ time; for compressed target length $N^{\beta+o(1)}$ it has exponent $1/\beta-o(1)$, not one uniform exponent as $\beta\to0$. Nothing in the hypothesis P=NP by itself licenses a PH verifier for this relation; asserting the displayed relation is in PH would imply $\oplus P\subseteq PH$.

\section*{SELECTOR-COVER-AND-TRAP12}
Let source promise $(Y,N_0)$, target promise $(Y',N')$, and seed-indexed maps $J_r$. Put $A_r=J_r^{-1}(Y')$, $B_r=J_r^{-1}(N')$.
\begin{theorem}
There exists an unbounded, source-adaptive selector $\sigma(x)$ mapping every promised source to the correct target promise side if and only if
\[
Y\subseteq\bigcup_r A_r,\qquad N_0\subseteq\bigcup_r B_r.
\]
If the target is total, $B_r=(Y\cup N_0)\setminus A_r$ and the second condition is $N_0\cap\bigcap_r A_r=\varnothing$.
\end{theorem}
\noindent Proof. Necessity: use the selected seed. Sufficiency: select one good seed per source, possibly via unbounded computation. QED. This does not construct an efficient selector; a toy system with one constant-YES and one constant-NO seed satisfies the two cover conditions for any source language, but choosing the correct seed computes the original source label. This toy is not a valid HASHIMAGE implementation.

\section*{RANDOM-INVENTORY-NO-TRAP12}
Let $N=2^n$, let $W_w$ denote the weight-$w$ truth tables ($L=\binom Nw$), and let $Y_t$ contain $W_w$. Let $A_1,\ldots,A_R$ have uniform marginal distributions over binary $b\times N$ matrices, independently of the input $f$.
\begin{theorem}
For any fixed $f\notin W_w$, with $\delta=(2^b-1)/(L+2^b-1)$,
\[
\Pr[\exists r:A_rf\notin A_r(Y_t)]\le R\delta.
\]
Hence with probability at least $1-R\delta$ a fixed genuine source NO $f$ is accepted under every public seed, defeating every selector restricted to that random inventory, even with unlimited computation.
\end{theorem}
\noindent Proof. For $g\in W_w$, the vectors $f+g$ are distinct and nonzero. For random $A$, the collision indicators $\mathbf1[A(f+g)=0]$ are pairwise independent, each with mean $2^{-b}$. Their sum $X$ has mean $L2^{-b}$ and variance $L2^{-b}(1-2^{-b})$. Cantelli gives $\Pr[X=0]\le(2^b-1)/(L+2^b-1)$. Union over $R$ seeds gives the claim. QED. For any NO distribution $\nu$, Fubini and Markov imply a random inventory traps a $1-\eta$ fraction of NO mass with probability at least $1-R\delta/\eta$.

\noindent With $t=N^\tau$ and $w=\lfloor(t-n+1)/n\rfloor$, all $W_w$ functions have size $\le t$ and $\log_2 L\ge(1-\tau-o(1))t$. If $b\le(1-\tau-\varepsilon)t$ and $R\le2^{\varepsilon t/2}$, the all-seed trap probability is at least $1-2^{-\varepsilon t/2+o(t)}$. Unlike full-image saturation, which uses $L\gg2^{2b}$, this fixed-source claim needs only $L\gg R2^b$. It repackages the prior CLOCKCOLLISION8/CLOCKEXCEPT10 concentration into the seed-portfolio quantifier; does not cover efficient short-seed or source-adaptive hash families. Explicitly specifying $R$ random matrices costs $\Theta(RNb)$ public bits, so do not claim a shrinking uniform reduction.

\section*{TESTS AND OUTCOME}
Runner local: 2026-10-18_clock_parity_inventory12_audit.py; raw JSON local: 2026-10-18_CLOCKPARITY12_raw.json; seed 202610182112. Exhaustive GF(4) tests: all 65,812 truth tables on $k=1,2,3,4$ variables, zero parity-reduction failures; four of four field labels admit sparse sections. Sandwich condition: 4,096 tested paired acceptance families on an eight-element domain, 729 feasible, zero discrepancies. With $N=5,b=2,w=2,f=7$, exact enumeration of all 1,024 matrices produced 93/1,024 exceptional seeds, below the $3/13$ Cantelli ceiling. In 10,000 randomized three-matrix inventories, 2,501 had any exceptional seed, below the $9/13$ union bound. Tests are finite, not runtime lower bounds.

\section*{STATUS AND NEXT TARGET}
PROVED/EXACT under declared hypotheses, independent review pending: supplied-circuit fingerprint equality counting-completeness; exact total/promise selector-cover characterization; all-seed random-inventory NO trap. RETIRED: interpreting $\exists C$ image membership as the same complexity problem as verifying a chosen $C$, using coverage existence as an efficient selector, treating independent full random matrices as a free short-index one-tape compiler. OPEN: one jointly NP/PH-low stage-uniform image verifier, source-aware YES/NO preservation, full one-work-tape fusion and the cofinal K lower-bound bridge. NEXT: SYMBOLIC-IMAGE-VERIFIER13. Build an explicitly source-aware, symbolically checkable, nonlinear quotient with easy witness evaluation but no cheap global sections; test against $\oplus P$, all-seed trapping, source-as-seed, Horn/2SAT/affine and alternative decoder escapes. P versus NP remains OPEN.

\begin{thebibliography}{4}
\bibitem{PZ} C. Papadimitriou and S. Zachos, Two Remarks on the Power of Counting, 1983.
\bibitem{VV} L. Valiant and V. Vazirani, NP is as easy as detecting unique solutions, Theor. Comput. Sci. 47 (1986), 85--93.
\bibitem{CW} J. Carter and M. Wegman, Universal Classes of Hash Functions, J. Comput. Syst. Sci. 18 (1979), 143--154.
\bibitem{CHMY} M. Cheraghchi et al., One-Tape Turing Machine and Branching Program Lower Bounds for MCSP, STACS 2021.
\end{thebibliography}
\end{document}
% FINAL_VERIFICATION_SENTINEL::CLOCK-GATE::CLOCKPARITY12::2026-10-18_2112::P_VS_NP_OPEN